\section{My support}

\textbf{4a. What others did that helped me work more effectively.} Four people changed how I worked, and I want to name the specific mechanism for each because a generic list of thanks is a Merit-ceiling move.

\textbf{Edward} is the reason the Cube ever spoke to the Pi. Around wk15 I had spent most of a weekend trying to make \texttt{pymavlink} read from \texttt{/dev/ttyAMA0} on Python 3.13, and it was Edward --- after I had filed three dead-end notes in Teams --- who pointed out that \texttt{pyserial} reads are broken on Python 3.13 and that the workaround was a \texttt{mavproxy} UDP bridge. That 40-minute conversation unlocked roughly two weeks of Pi testing. Edward's throwaway remark that ``just because the test passes on laptop doesn't mean the hardware does'' reframed how I thought about simulation fidelity for the rest of the project --- I stopped treating SITL parity as evidence and started treating it as a hypothesis.

\textbf{Robin} helped me by refusing to integrate against my first state machine. In wk17 he pushed back until I could prove the extra states caused materially different decisions; his pushback is the reason the \texttt{main.py} refactor (1411\,$\to$\,754 lines, 6 modules) happened at all. What I had framed as a technical disagreement was actually a maintainability one, and I preferred the technical framing because it protected me from noticing I had built something only I could read. \textbf{Demetro} accepted drafts in my voice without turning collaboration into a status contest. \textbf{Our pilot} operated the browser ground station unaided during field-day dry runs. \textbf{Sid} delivered the single DJI flight video with SRT telemetry that anchored FOV calibration and the retrained model at mAP50\,=\,0.995.

\textbf{4b. What I did that helped others work more effectively.} I will put three specific feedback episodes on record, because ``I gave constructive feedback when needed'' is the single most common Merit-ceiling sentence in UK engineering reflection.

\textit{Feedback to Demetro --- FDR speaking scripts.} When Demetro asked for help with his two FDR slides I built the layouts and wrote full speaking scripts (246 + 219 words, 3.1\,min total, commits \texttt{b2a9cd4}, \texttt{ba667f2}), iterating 4--5 times. On the third pass he pushed back: the wording was not how he spoke. The observable behaviour change came on FDR day --- he delivered the two slides confidently, in his own cadence, and a week later prepared his own next set using specific numbers rather than adjectives (a habit I had modelled in draft three). What I learned about my own delivery is harder: my first three drafts were in my register, not his, and I now see that good feedback to a peer presenter means writing in their voice, not polishing yours. The scarce resource was register, not prose quality.

\textit{Feedback to Robin --- CENTERING timeout and the handoff contract.} In wk19 I noticed Robin's flight script had no timeout guard on the \texttt{CENTERING} state: if the target was lost mid-descent the loop would hang. I sent him a Teams message naming the exact variable (\texttt{consecutive\_lost}) and the exact line, and suggested a time-based timeout with fallback to ascend-and-retry. He fixed it the same afternoon. The behaviour change that mattered came a week later: he pulled me aside to say he had caught a similar timeout bug in a separate module because he was now ``looking for that shape of bug.'' The generalisable lesson is that when giving feedback to a high-performing peer, specificity is the scarce resource, not softness. I now see that my instinct to ``finish debugging before asking'' is not professionalism --- it is a cost I force teammates to pay for my own comfort with solitude, and I should be asking the moment I notice I have been stuck for twenty minutes, not the moment frustration has built up far enough that the help-request feels validated to me.

\textit{Feedback to the pilot --- ground-station button layout and jargon removal.} When I showed the pilot the first \texttt{pi\_flight.py} dashboard he said the button grid was too dense to hit reliably while holding an RC transmitter. I proposed a minimal four-button set (Y / N / M / L) mapped to the RC thumb-reach pattern, and wired an emergency RTL into the override guard so mode 9 (LAND) from the RC switch would be honoured. The subtler adaptation was lexical: my first draft used phrases like ``geofence repulsor active within 10\,m buffer'' which I rewrote as ``the drone pushes itself away from the red fence if it gets within 10\,m.'' The rewritten card was read aloud in one pass without a pause, and during the field-day dry runs the pilot triggered the abort three times without looking at the screen --- the moment my frame shifted from designing for me-as-operator to designing for a pilot who does not have to notice the UI.

The quieter half of this work was infrastructure for the team's memory rather than its code: \texttt{brain-dump.md}, \texttt{MEMORY.md}, the 500-line blueprint system, and \texttt{FIELD\_QUICK\_REF.md} for no-internet field use.

\textbf{What I leave believing.} I came as a Ukrainian-trained engineer whose reflex under failure was ``one person delivers''; I leave believing its near-inversion --- \emph{the engineer who delivers alone is one recruitment away from being the single point of failure in whatever she builds}. The wk17 FSM conflict and Robin's wk20 unaided handoff told me, in different registers, that the same act which delivered the report was the act that kept my teammates standing at the edge of their own project. Falsifiable internship signal: in my first six weeks, does at least one teammate modify a module I own without asking permission? If yes --- even once --- I have started to become an engineer whose work enables other people's work, which is the only kind of engineering I now think is worth aspiring to.
